\subsection{Directional Measurement and Estimator}\label{sec:adjustment-design}

The model's exposure ordering concerns a target's movement relative to other
firms, which we measure by holding their article distributions fixed while
the target's distribution changes.
This preserves a common origin for each comparison, separating the target's
approach from the reference firm's contemporaneous movement.
Let $\widehat C_{iy}$ be firm $i$'s annual distribution over unit-normalized
article vectors and keep every reference law $\widehat C_{j0}$ at the 2018
origin, and for reference $j\ne i$ define
\begin{equation}\label{eq:frozen-approach}
  D_{ijy}=W_2^2(\widehat C_{iy},\widehat C_{j0}),\qquad
  A_{ijy}=D_{ij,y-1}-D_{ijy}.
\end{equation}
Positive $A$ means that target $i$ approaches reference $j$, with $A$ measured
in squared embedding-distance units under the squared Euclidean ground cost.
Under the two-opportunity law \eqref{eq:axis-observation}, changes in
distribution shape reveal movement even though all centroids coincide.

News of a free cash flow shock arriving early in a year can precede more of
the target's coverage than news arriving near its end.
We reflect this timing by weighting each reference return by the fraction
of target articles still to appear.
For return date $s$ in destination year $y$, denote that fraction by
$h_{iys}$ and the reference's daily simple decimal excess return by
$r^e_{js}$, and set
\begin{equation}\label{eq:dated-financial}
  h_{iys}=\frac{1}{n_{iy}}\sum_{\ell=1}^{n_{iy}}
  \mathbf1\{d_{iy\ell}>s\},\qquad
  R_{ijy}=\sum_{s\in y}h_{iys}r^e_{js}.
\end{equation}
We sum each reference firm's daily excess returns, weighting them by the
fraction of the target firm's articles published later in the year.
This is an additive return measure, not a compounded annual return.
The strict inequality excludes same-day news without assigning an intraday
arrival time.
Because dates remain paired with articles, each redraw changes both the
information distribution and its article-timing weights.

Centering subtracts the cross-reference mean within each target-year,
so the coefficient compares which reference profiles the target approaches
more closely with which references have higher weighted excess returns.
References receive equal weight within each target-year.
For $V=A,R$, write $\bar V_{iy}=(N-1)^{-1}\sum_{j\ne i}V_{ijy}$ and
$\widetilde V_{ijy}=V_{ijy}-\bar V_{iy}$.
With $Y$ transitions, the annual moments and pooled estimator are
\begin{equation}\label{eq:directional-ratio}
  q_y=\frac{\sum_i\sum_{j\ne i}\widetilde R_{ijy}\widetilde A_{ijy}}{N(N-1)},\quad
  d_y=\frac{\sum_i\sum_{j\ne i}\widetilde R_{ijy}^2}{N(N-1)},\quad
  \widehat\beta=\frac{\sum_y q_y}{\sum_y d_y}.
\end{equation}
The coefficient solves the aggregate centered moment $\sum_y(q_y-\beta
d_y)=0$ and minimizes the pooled squared residual after target-year intercepts.
A return date with zero article-timing weight contributes zero to the
weighted excess-return measure without invalidating the comparison.
A zero pooled financial denominator makes the coefficient undefined.
When annual ratios differ, pooling weights years by $d_y$, giving reference
firms with greater relative financial dispersion greater denominator
influence despite equal reference weights.

Each firm enters many comparisons, both as a target and as a reference.
Article redraws therefore reuse each sampled firm-year distribution wherever
that firm appears, preserving the resulting dependence.
These redraws assess article-selection variability for the observed firms
and financial histories under the scheme in Section~\ref{sec:data}.

Uncertainty under independently sampled complete firm histories concerns a
different target.
Appendix~\ref{app:annual-gmm} uses a complete-triple representation to
examine a variance approximation under independently sampled complete firm
histories, retaining dependence across years and both roles.
The simulation evaluates the two procedures against their respective known
targets.

\subsection{Restricted-Form Hypothesis}\label{sec:hypothesis}

In the restricted two-opportunity economy, when one opportunity improves
relative to the other, firms with greater initial exposure to it have more
favorable return responses, while other firms expanding toward it approach those
firms' initial news distributions.
Because Proposition~\ref{prop:directional-alignment} signs the moments that
\eqref{eq:directional-ratio} pools, its conditions imply a nonnegative population
ratio of expected pooled moments whenever the denominator is positive.
In the simulated economy, where this coefficient is strictly positive in every
history under common reporting, we therefore state the hypothesis as
$H_0:\beta\leq0$ against $H_1:\beta>0$.

Since changes in reporting alone can also produce a positive coefficient at
fixed capacities, we assess the estimator against a positive operating target
under common reporting, a zero target under stationary reporting, and a positive
reported target without operating movement.
In the corporate archive, where the proposition's conditions are not verified
and no established design samples independent firm histories, we report the
estimate with its sensitivity to article selection
rather than test $H_0$.
